\section{Introduction}

In the \emph{$0$-Extension} problem (\ze), we are given an undirected edge-weighted graph $G=(V,E,c)$, a set $T\subseteq V$ of its vertices called \emph{terminals}, and a metric $D$ on terminals, and the goal is to find a mapping $f: V\to T$ that maps each vertex to a terminal in $T$, such that each terminal is mapped to itself (i.e., $f(t)=t$ for all $t\in T$), and the sum $\sum_{(u,v)\in E}c{(u,v)}\cdot D(f(u),f(v))$ is minimized.

The \ze problem was first introduced by Karzanov \cite{karzanov1998minimum}. It is a generalization of the \emph{multi-way cut} problem (by setting $D(t,t')=1$ for all pairs $t,t'\in T$) \cite{dahlhaus1994complexity,cualinescu1998improved,freund2000lower,buchbinder2013simplex,angelidakis2017improved,buchbinder2017simplex,berczi2020improving}, and a special case of the \emph{metric labeling} problem \cite{kleinberg2002approximation,chekuri2004linear,archer2004approximate,karloff2006earthmover,chuzhoy2007hardness}.
C\u{a}linescu, Karloff and Rabani \cite{calinescu2005approximation} gave the first approximation algorithm for \ze, achieving a ratio of $O(\log |T|)$, by rounding the solution of a \emph{semi-metric LP relaxation} (\textsf{LP-Metric}), which is presented below.
%
\begin{eqnarray*}
	\mbox{(\textsf{LP-Metric})}\quad &  \text{minimize}\quad  \sum_{(u,v)\in E}c(u,v)\cdot \delta(u,v)\\
	& s.t. \quad (V,\delta) \text{ is a semi-metric space}\\
	& \delta(t,t')=D(t,t'), \quad\forall t,t'\in T
\end{eqnarray*}
 %
Fakcharoenphol, Harrelson, Rao and
Talwar \cite{fakcharoenphol2003improved} later gave a modified rounding algorithm on the same LP, improving the ratio to $O(\log |T|/\log\log |T|)$, which is the current best-known approximation. On the other hand, this LP was shown to have integrality gap $\Omega(\sqrt{\log |T|})$ \cite{calinescu2005approximation}, and this was recently improved to $\Omega((\log |T|)^{2/3})$ by Schwartz and Tur \cite{schwartz2021metric}.
Another LP relaxation called \emph{earthmover distance relaxation} (\textsf{LP-EMD}) was considered by Chekuri, Khanna, Naor and Zosin \cite{chekuri2004linear} and utilized to obtain an $O(\log |T|)$-approximation of the metric labeling problem (and therefore also the \ze problem). It has been shown \cite{karloff2009earthmover} by Karloff, Khot, Mehta and Rabani that this LP relaxation has an integrality gap $\Omega(\sqrt{\log |T|})$. Manokaran, Naor, Raghavendra
and Schwartz \cite{manokaran2008sdp} showed that the integrality gap of this LP relaxation leads to the same hardness of approximation result, assuming the Unique Game Conjecture.


In addition to being an important problem on its own, the \ze problem and its two LP relaxations are also closely related to the construction of \emph{cut/flow vertex sparsifiers}, a central problem in the paradigm of graph compression.
Given a graph $G$ and a set $T\subseteq V(G)$ of terminals, a cut sparsifier of $G$ with respect to $T$ is a graph $H$ with $V(H)=T$, such that for every partition $(T_1,T_2)$ of $T$, the size of the minimum cut separating $T_1$ from $T_2$ in $G$ and the size of the minimum cut separating $T_1$ from $T_2$ in $H$, are within some small factor $q$, which is also called the \emph{quality} of the sparsifier\footnote{flow sparsifiers has a slightly more technical definition, which can be found in \cite{hagerup1998characterizing,leighton2010extensions,chuzhoy2012vertex,andoni2014towards}.}.
Moitra \cite{moitra2009approximation} first showed that every graph with $k$ terminals admits a cut sparsifier with quality bounded by the \emph{integrality gap} of its \textsf{LP-Metric} (hence $O(\log k/\log\log k)$). Later on, Leighton and Moitra \cite{leighton2010extensions}, and Makarychev and Makarychev \cite{makarychev2010metric} concurrently obtained the same results for flow sparsifiers, and then Charikar, Leighton, Li and Moitra \cite{charikar2010vertex} showed that the best flow sparsifiers can be computed by solving an LP similar to \textsf{LP-EMD}. On the lower bound side, it was shown after a line of work \cite{leighton2010extensions,charikar2010vertex,makarychev2010metric} that there exist graphs with $k$ terminals whose best flow sparsifier has quality $\tilde\Omega(\sqrt{\log k})$.


A major open question on vertex sparsifiers is: 

\vspace{-10pt}
\[\emph{\underline{Q1.} Can better quality sparsifiers be achieved by allowing a small number of Steiner vertices}? 
\]
\vspace{-10pt}

In other words, what if we no longer require that the sparsifier $H$ only contain terminals, but just require that $H$ contain all terminals and its size be bounded by some function $f$ on the number of terminals (for example, $f(k)=2k,k^2$ or even $2^k$)? 
Chuzhoy \cite{chuzhoy2012vertex} constructed $O(1)$-quality cut/flow sparsifiers with size dependent on the number of terminal-incident edges in $G$.
Andoni, Gupta and Krauthgamer \cite{andoni2014towards} showed the construction for $(1+\eps)$-quality flow sparsifiers for quasi-bipartite graphs. For general graphs, they constructed a sketch of size $f(k,\eps)$ that stores all feasible multicommodity flows up to a factor of $(1+\eps)$, raising the hope for a special type of $(1+\eps)$-quality flow sparsifier, called contraction-based flow sparsifiers, of size $f(k,\eps)$ for general graphs, which was recently invalidated by Chen and Tan \cite{chen20241+}, who showed that contraction-based flow sparsifiers whose size are bounded by any function $f(k)$ must have quality $1+\Omega(1)$. But it is still possible for such flow sparsifiers with constant quality and finite size to exist. Prior to this work, Krauthgamer and Mosenzon \cite{krauthgamer2023exact} showed that there exist $6$-terminal graphs $G$ whose quality-$1$ flow sparsifiers must have an arbitrarily large size.

Given the concrete connection between the \ze problem and  cut/flow sparsifiers, it is natural to ask a similar question for \ze: 

\vspace{-10pt}
\[\emph{\underline{Q2.} Can better approximation of \textnormal{\ze} be achieved by allowing a small number of Steiner vertices?}\]
\vspace{-10pt}

In this paper, we formulate and study the following variant of the \ze problem, which we call the \emph{0-Extension with Steiner Nodes} problem ($\zesn$). (We note that a similar variant was mentioned in \cite{andoni2014towards}, and we provide a comparison between them in more detail in \Cref{sec: compare}.)
Upfront we are also given a function $f: \mathbb{Z}\to \mathbb{Z}$ with $f(k)\ge k$ holds for all $k\in \mathbb{Z}$, specifying the number of Steiner vertices allowed.

\vspace{-10pt}

\paragraph{0-Extension with Steiner Nodes.}
In an instance of $\zesn(f)$, the input consists of an edge-weighted graph $G=(V,E,c)$, a subset $T\subseteq V$ of $k$ vertices, that we call \emph{terminals}, and a metric $D$ on terminals in $T$, which is exactly the same as \ze.
%We denote by $\dist_{\ell}(\cdot,\cdot)$ the shortest-path (in $G$) distance metric on $V$ induced by the lengths $\{\ell_e\}_{e\in E}$.
A solution to the instance $(G,T,D)$ consists of 
\begin{itemize}
\item a partition $\fset$ of $V$ with $|\fset|\le f(k)$, such that distinct terminals of $T$ belong to different sets in $\fset$; we call sets in $\fset$ \emph{clusters}, and for each vertex $u\in V$, we denote by $F(u)$ the cluster in $\fset$ that contains it;
\item a semi-metric $\delta$ on the clusters in $\fset$, such that for each pair $t,t'\in T$, $\delta(F(t),F(t'))= D(t,t')$.
\end{itemize}
%
We define the \emph{cost} of a solution $(\fset,\delta)$ as $\vol(\fset,\delta)=\sum_{(u,v)\in E}c(u,v)\cdot\delta(F(u),F(v))$, and its \emph{size} as $|\fset|$.
The goal is to compute a solution $(\fset,\delta)$ with size at most $f(k)$ and minimum cost. 

The difference between $\zesn(f)$ and \ze is that, instead of enforcing every vertex to be mapped to a terminal, in $\zesn(f)$ we allow vertices to be mapped $(f(k)-k)$ non-terminals (or Steiner nodes), which are the clusters in $\fset$ that do not contain terminals. 
We are also allowed to manipulate the distances between these non-terminals, conditioned on not destroying the induced metric $D$ on terminals. Clearly, when $f(k)=k$, the $\zesn(f)$ problem degenerates to the \ze problem.
%This additional freedom makes it possible to get around the $\Omega(\sqrt{\log k})$ integrality gap

%We still require that a metric be defined on these clusters, such that the induced metric on terminal clusters (clusters in $\fset$ that contain terminals) is exactly $D$. 

It is easy to see that (\textsf{LP-Metric}) is still an LP relaxation for $\zesn(f)$, as each solution $(\fset,\delta)$ to $\zesn$ naturally corresponds to a semi-metric $\delta'$ on $V$ (where we can set $\delta'(u,u')=\delta(F(u),F(u'))$ for all pairs $u,u'\in V$). 
Denote by $\textsf{IG}_f(k)$ the worst integrality gap for (\textsf{LP-Metric}) to any $\zesn(f)$ instance with at most $k$ terminals.
In fact, similar to the connection between the integrality gap of (\textsf{LP-Metric}) and the quality achievable by flow sparsifiers \cite{moitra2009approximation,leighton2010extensions},
%following from the results in \cite{andoni2014towards}\footnote{This follows from LP(2) and Proposition 4.2 of \cite{andoni2014towards}. See \Cref{sec: compare} for a detailed discussion.}, 
it was recently shown by Chen and Tan \cite{chen20241+} that the value of $\textsf{IG}_f(k)$ is also closely related to the quality achievable by flow sparsifiers with Steiner nodes.
Specifically, for any function $f$, every graph $G$ with $k$ terminals has a quality-$\big((1+\eps)\cdot \textsf{IG}_f(k)\big)$ flow sparsifier with size bounded by $(f(k))^{(\log k/\eps)^{k^2}}$. 
This means that any positive answer to question Q2 (by proving that $\textsf{IG}_f(k)=o(\log k/\log\log k)$ for some $f$) also gives a positive answer to question Q1.

This makes it tempting to study the $\zesn$ problem. Specifically, can we prove any better-than-$O(\log k/\log\log k)$ upper bound for $\textsf{IG}_f(k)$, for any function $f$?
To the best of our knowledge, no such bound is known for any $f$, leaving the problem wide open. 
%Therefore, it is natural to ask if the additional freedom in solutions can significantly improve the old $\Omega(\sqrt{\log k})$ integrality gap. Specifically, for some small function $f$ such as $f(k)=O(k)$, can we hope to improve the integrality gap of (\textsf{LP-Metric}) to $\zesn(f)$ to $O(1)$?
In fact, no non-trivial lower bound on $\textsf{IG}_f(k)$ is known for even very small function like $f(k)=O(k)$.


\subsection{Our Results}

In this paper, we make a first step in investigating the value of  $\textsf{IG}_f(k)$, by giving a superconstant lower bound on $\textsf{IG}_f(k)$ for near-linear functions $f$.
Our main result is summarized in the following theorem.
%\Cref{thm: 0-ext with Steiner lower}

\begin{theorem}
\label{thm: 0-ext with Steiner lower}
For any $0<\eps<1$ and any size function $f:\mathbb{Z}^+\to \mathbb{Z}^+$ with $f(k)=O(k\log^{1-\eps} k)$, 
the integrality gap of the LP-relaxation (\textnormal{\textsf{LP-Metric}}) is $\textnormal{\textsf{IG}}_f(k)=\Omega(\eps \log\log k)$.
%for every sufficiently large $k$, there exists an instance $(G,T,D)$ of $\zesn$ with $|T|=k$ and size bound $f(k)$, such that the integrality gap of  (\textnormal{\textsf{LP-Metric}}) is $\Omega(\eps \log\log k)$.
\end{theorem}

%To the best of our knowledge, there is no known upper bound better than $O(\log k/\log\log k)$ for any function $f$.

We remark that our lower bound for the integrality gap of $\zesn(f)$ does not imply a size lower bound for $O(\log\log k)$-quality flow sparsifiers. However, if the same lower bound can be proved for a slightly generalized version of $\zesn(f)$, that was proposed in \cite{andoni2014towards} and analyzed in \cite{chen20241+}, then it will imply an $\Omega(k\log^{1-\eps} k)$ size lower bound for $O(\log\log k)$-quality flow sparsifiers. We provide a detailed discussion in \Cref{sec: compare}.

\subsection{Technical Overview}

We now discuss some high-level ideas in the proof of  \Cref{thm: 0-ext with Steiner lower}. Given any $k$, we will construct an unweighted graph $G$ on $n$ vertices (where $n\approx k\log k$) and $k$ terminals, and show that any solution of the $\zesn$ problem with size $O(k\log^{1-\eps} k)$ has cost lower bounded by $\Omega(\log\log k)$ times the number of edges in $G$.

Our hard instance is a constant degree expander (with an arbitrary set of its $k$ vertices as terminals). There are two main reasons to choose such a graph. First, in the previous work \cite{calinescu2005approximation} for proving the $\Omega(\sqrt{\log k})$ integrality gap lower bound for the $0$-Extension problem, an graph called ``expander with tails'' has been used. Though the tails in their construction appear useless for our purpose, as we allow a super-linear number of Steiner vertices which easily accomodate a single-edge tail for each terminal, the expander graph turns out to be still the core structure that is hard to compress.
Second, in another previous work \cite{archer2004approximate} it was shown that $0$-Extension problem on minor-free graphs has integrality gap $O(1)$, so our hard example has to contain large cliques as minor, which makes expanders, the graphs with richest structure,  a favorable choice. For technical reasons, we need some additional properties like Hamiltonian and high-girth. See \Cref{subsec: instance} for more details.

Next we want to lower bound the cost of an $\zesn$ solution of size $O(k\log^{1-\eps} k)$. This is done in two steps. We first consider a special type of solutions that really ``comes from the graph''. Recall that a solution consists of a partition $\fset$ of $V(G)$ into clusters and a metric $\delta$ on clusters in $\fset$. Specifically, in this special type of solutions, we require that each cluster $F\in \fset$ corresponds to a distinct vertex $v_F$ (called its \emph{center}) in $G$, and for all pairs $F,F'$ the metric $\delta(F,F')$ coincides with the shortest-path distance between $v_F, v_{F'}$ in $G$.
In the first step, we show that all special solutions have cost $\Omega(n\cdot \eps \log\log k)$. Intuitively, as the graph is a constant degree expander, every center $v_F$ is within distance $(\eps/100)\cdot \log\log k$ to at most $(\log k)^{\eps/10}$ other centers, but its cluster $F$ contains $(\log k)^{\eps}$ vertices on average and so it has $\Omega((\log k)^{\eps})$ inter-cluster edges. Therefore, most edges will have new length at least $(\eps/100)\cdot \log\log k$ (in $\delta$), making the total new edge length $\Omega(n \cdot \eps \log\log k)$ (as the number of inter-cluster edges in a balanced expander partition is $\Omega(n)$). Careful calculations are needed to turn these informal arguments into a rigorous proof. See \Cref{subsec: canonical} for more details.

In the second step, we show that the general solutions can actually be reduced to the special type of solutions considered in the first step, losing only an $O(1)$ factor in its cost. In fact, it has been recently shown in \cite{chen20241+} that, to analyze the cost of any $\zesn$ instance, it suffices to consider $\zesn$ solutions whose metric $\delta$ is embeddable into a geodesic structure of the terminal-induced shortest path distance metric called \emph{tight span}. Our main contribution here, on a conceptual level, is showing that \emph{for a graph with high girth, its tight span structure locally coincides with the graph structure itself}. In this sense, we can compare any solution to some special solution considered in Step 1. When the picture is lifted to the global level, in such a comparison it turns out that we will loss a factor approximately the \emph{diameter-girth ratio}, which we can manage to get $O(1)$ with an additional short-cycle-removing step in the construction of the expander. We believe this diameter-girth ratio quantifies how ``local'' a graph structure is, and should be of independent interest to other graph problems on shortest-path distances.
To carry out the technical steps, we employ a notion called \emph{continuazation} of a graph recently studied in \cite{chen2023towards}.
See \Cref{subsec: mapping} for more details.


\iffalse{old version}

In the 0-Extension with Steiner Nodes problem, the input consists of 
\begin{itemize}
	\item a graph $G=(V,E)$ on $|V|=n$ vertices; 
	\item a subset $T\subseteq V$ of $k$ vertices, that we call \emph{terminals}; and
	\item lengths $\{\ell_e\}_{e\in E}$ of edges of $G$. 
\end{itemize}

We denote by $\dist_{\ell}(\cdot,\cdot)$ the shortest-path (in $G$) distance metric on $V$ induced by the lengths $\{\ell_e\}_{e\in E}$.

A solution to the $\zesn$ problem consists of (i) a partition $\fset$ of the vertex set $V$, such that distinct terminals of $T$ belong to different sets in $\fset$; we call sets in $\fset$ \emph{clusters}, and for each vertex $u\in V$, we denote by $F(u)$ the cluster in $\fset$ that contains it; and (ii) a semi-metric $\delta$ on the clusters in $\fset$, such that for each pair $t,t'\in T$, 
%if we denote by $F(t)$ the cluster that contains $t$, and by $F(t')$ the cluster that contains $t'$, then 
$\delta(F(t),F(t'))= \dist_{\ell}(t,t')$.


We define the \emph{cost} of a solution $(\fset,\delta)$ as $\vol(\fset,\delta)=\sum_{(u,v)\in E}\delta(F(u),F(v))$, and the \emph{size} of a solution $(\fset,\delta)$ to be $|\fset|$.
The goal is to compute a solution $(\fset,\delta)$ with small size and cost. Ideally, we would like that $|\fset|$ only depends on $k$, and that $\vol(\fset,\delta)$ is only a constant factor away from the value $\sum_{(u,v)\in E}\ell_{(u,v)}$, which can be viewed as the cost of the original graph.
We say that a solution $(\fset,\delta)$ achieves \emph{stretch} $\rho$ iff
$\vol(\fset,\delta)
\le \rho\cdot
\sum_{(u,v)\in E}\ell_{(u,v)}$.

In an extreme setting, if we allow $|\fset|=n$, then we can achieve stretch $1$ by letting $\fset=\set{\set{v}\mid v\in V}$ and $\delta(\set{u},\set{v})= \dist_{\ell}(u,v)$ for all pairs $u,v\in V$. In the other extreme setting, it has been shown by \cite{fakcharoenphol2003improved} that a solution $(\fset,\delta)$ with size $|\fset|=k$ and stretch $O(\log k/\log\log k)$ always exist and can be found efficiently.In this note, we prove the following theorem.

\fi
